En principe, le processus de transport doit conserver l'energie, donc
l'impulsion $p$, et le moment magnetique $\mu$ de la particule,
ces deux quantites etant des invariants du mouvement. La conservation
du moment magnetique impose la relation suivante, \[
\mu=\left.\frac{1-\xi^{2}}{B_{\min}}\right|_{l=l_{1}}=\left.\frac{1-\xi^{2}}{B_{\min}}\right|_{l=l_{2}}\]
 ou $l_{1}$ et $l_{2}$ sont les indices correspondant aux deux positions
radiales entre lesquelles l'electron se deplace. La valeur de $\xi$
est donc une fonction de la position radiale $r$, a travers la dependance
du champ magnetique $B$. Ainsi, \[
\xi_{j_{l_{2}}}=\sigma_{j_{l_{1}}}\sqrt{1-\Psi_{l_{1},l_{2}}\left(1-\xi_{j_{l_{1}}}^{2}\right)}\]
 ou \[
\Psi_{l_{1},l_{2}}=\frac{B_{\min}\left(r_{l_{2}}\right)}{B_{\min,}\left(r_{l_{1}}\right)}\]
 et \[
\sigma_{j_{l_{1}}}=sgn\left(\xi_{j_{l_{1}}}\right)\]


De proche en proche, un electron circulant au centre du plasma peut
ainsi devenir un electron piege au bord de celui-ci. Cependant, les
corrections liees a la conservation du moment magnetique concernent
essentiellement les electrons ayant une grande energie perpendiculaire,
soit $\xi\approx0.$ En effet, c'est dans cette region uniquement
que la variation relative de $\xi$ d'un rayon a l'autre est la plus
importante. Pour les electrons les plus rapides, donc $\left|\xi\right|\sim1,$
cette correction est faible, voire negligeable, et le modele simple
du paragraphe precedent est alors largement suffisant. Lorsque l'electron
migre du centre vers le bord du plasma, $\xi$ tend a augmenter, donc
son impulsion parallele $p_{\Vert}.$ Dans le cadre du modele simple
du paragraphe precedent, on sous-estime l'elargissement du courant
toro\"{\i}dal qui en resulte au bord du plasma.

Sous forme generale, l'equation conservative du transport radial moyennee
sur les trajectoires vaut alors \[
\frac{\partial\left(rp^{2}f_{0}^{\left(0\right)}\right)}{\partial t}+p^{2}\frac{\partial\left(rS_{r}^{\left(0\right)}\right)}{\partial r}=0\]
 ou \[
S_{r}^{\left(0\right)}\left(p,\xi\left(r\right),r\right)=-D_{rr}^{\left(0\right)}\left(p,\xi\left(r\right),r\right)\frac{\partial f_{0}^{\left(0\right)}\left(p,\xi\left(r\right),r\right)}{\partial r}+F_{r}^{\left(0\right)}\left(p,\xi\left(r\right),r\right)f_{0}^{\left(0\right)}\left(p,\xi\left(r\right),r\right)\]


On en deduit ainsi la forme discretisee sur la grille,

\[
r_{l+1/2}p_{i+1/2}^{2}\frac{f_{0,l+1/2,j_{l+1/2}+1/2}^{\left(0\right)(k+1)}-f_{0,l+1/2,j_{l+1/2}+1/2}^{\left(0\right)(k)}}{\Delta t}=-p_{i+1/2}^{2}\frac{r_{l+1}\left.S_{r}^{\left(0\right)}\right|_{l+1,j_{l+1}+1/2}^{(k+1)}-r_{l}\left.S_{r}^{\left(0\right)}\right|_{l,j_{l}+1/2}^{(k+1)}}{\Delta r_{l+1/2}}\]
 avec \begin{eqnarray*}
\left.S_{r}^{\left(0\right)}\right|_{l+1,j_{l+1}+1/2}^{(k+1)} & = & -D_{rr,l+1,j_{l+1}+1/2}^{\left(0\right)}\frac{f_{0,l+3/2,j_{l+3/2}+1/2}^{\left(0\right)(k+1)}-f_{0,l+1/2,j_{l+1/2}+1/2}^{\left(0\right)(k+1)}}{\Delta r_{l+1}}+F_{r,l,j_{l+1}+1/2}^{\left(0\right)}f_{0,l+1,j_{l+1}+1/2}^{\left(0\right)(k+1)}\\
\left.S_{r}^{\left(0\right)}\right|_{l,j_{l}+1/2}^{(k+1)} & = & -D_{rr,l,j_{l}+1/2}^{\left(0\right)}\frac{f_{0,l+1/2,j_{l+1/2}+1/2}^{\left(0\right)(k+1)}-f_{0,l-1/2,j_{l-1/2}+1/2}^{\left(0\right)(k+1)}}{\Delta r_{l}}+F_{r,l,j_{l}+1/2}^{\left(0\right)}f_{0,l,,j_{l}+1/2}^{\left(0\right)(k+1)}\end{eqnarray*}
 ou \[
\left\{ \begin{array}{l}
\xi_{j_{l+3/2}+1/2}=\sigma_{j_{l+1/2}+1/2}\sqrt{1-\Psi_{l+1/2,l+3/2}\left(1-\xi_{j_{l+1/2}+1/2}^{2}\right)}\\
\xi_{j_{l+1}+1/2}=\sigma_{j_{l+1/2}+1/2}\sqrt{1-\Psi_{l+1/2,l+1}\left(1-\xi_{j_{l+1/2}+1/2}^{2}\right)}\\
\xi_{j_{l}+1/2}=\sigma_{j_{l+1/2}+1/2}\sqrt{1-\Psi_{l+1/2,l}\left(1-\xi_{j_{l+1/2}+1/2}^{2}\right)}\\
\xi_{j_{l}-1/2}=\sigma_{j_{l+1/2}+1/2}\sqrt{1-\Psi_{l+1/2,l-1/2}\left(1-\xi_{j_{l+1/2}+1/2}^{2}\right)}\end{array}\right.\]


A la limite $\Psi=1,$ c'est-a-dire en negligeant la variation locale
du champ magnetique, on retrouve bien le cas simple precedent.

De maniere generale, il n'est pas possible d'effectuer un developpement
de Taylor autour de la valeur en $j_{l+1/2}+1/2$ pour exprimer les
termes en $j_{l-1/2}+1/2$, $j_{l}+1/2$, $j_{l+1}+1/2$ et $j_{l+3/2}+1/2,$
car les grilles peuvent etre telles qu'une variation radiale $\Delta r_{l+1/2}$
corresponde a une changement important de la distribution en $\xi$.
L'interpolation entre les grilles vaut alors \[
\left\{ \begin{array}{l}
f_{0,l+1,j_{l+1}+1/2}^{\left(0\right)(\infty)}=\left(\frac{\Delta r_{l+1/2}}{\Delta r_{l+1/2}+\Delta r_{l+3/2}}\right)f_{0,l+3/2,j_{l+1}+1/2}^{\left(0\right)(\infty)}+\left(\frac{\Delta r_{l+3/2}}{\Delta r_{l+1/2}+\Delta r_{l+3/2}}\right)f_{0,l+1/2,j_{l+1}+1/2}^{\left(0\right)(\infty)}\\
f_{0,l,,j_{l}+1/2}^{\left(0\right)(\infty)}=\left(\frac{\Delta r_{l-1/2}}{\Delta r_{l-1/2}+\Delta r_{l+1/2}}\right)f_{0,l+1/2,j_{l}+1/2}^{\left(0\right)(\infty)}+\left(\frac{\Delta r_{l+1/2}}{\Delta r_{l-1/2}+\Delta r_{l+1/2}}\right)f_{0,l-1/2,j_{l}+1/2}^{\left(0\right)(\infty)}\end{array}\right.\]


d'ou

\[
r_{l+1/2}p_{i+1/2}^{2}\frac{f_{l+1/2,j_{l+1/2}+1/2}^{\left(0\right)(k+1)}-f_{l+1/2,j_{l+1/2}+1/2}^{\left(0\right)\left(k\right)}}{\Delta t}+p_{i+1/2}^{2}\sum_{l^{\prime}=l-1}^{l^{\prime}=l+1}\overline{\overline{M}}_{r,l^{\prime}+1/2}^{\left(0\right)}f_{0,l^{\prime}+1/2,j_{l^{\prime}+1/2}+1/2}^{\left(0\right)(k+1)}=0\]
 les coefficients de matrice $\overline{\overline{M}}_{r}^{\left(0\right)}$
valant dans ce cas,{*}{*}{*}{*}{*}{*}{*}{*}{*}{*}{*}{*}{*}{*} A ffaire
{*}{*}{*}{*}{*}{*}{*}{*}{*}{*}{*}{*}{*}{*}{*}{*}{*}
